% !PRINTED 230
\weberpage{230}
%<<CONT
60\textsuperscript{sten} Grades von fünf Ziffern. Diese Permutationsgruppe ergiebt sich, wenn wir mit den Nebengruppen
\[
Q, \quad Q\Theta, \quad Q\Theta^2, \quad Q\Theta^3, \quad Q\Theta^4
\]
die sämmtlichen Elemente $\sigma$ der Ikosa\"edergruppe verbinden, also
\[
Q\sigma, \quad Q\Theta\sigma, \quad Q\Theta^2\sigma, \quad Q\Theta^3\sigma, \quad Q\Theta^4\sigma
\]
bilden, und die dadurch bewirkte Permutation dieser Nebengruppen untersuchen. Für die erzeugenden Substitutionen der Ikosa\"edergruppe $\sigma = \Theta, \psi, \chi$ erhalten wir so die Permutationen
\[
(1, \ \Theta, \ \Theta^2, \ \Theta^3, \ \Theta^4) \qquad \sigma = \Theta
\]
\[
\begin{pmatrix} Q, & Q\Theta, & Q\Theta^2, & Q\Theta^3, & Q\Theta^4 \\ Q, & Q\Theta^4, & Q\Theta^3, & Q\Theta^2, & Q\Theta \end{pmatrix} \qquad \sigma = \psi
\]
\[
\begin{pmatrix} Q, & Q\Theta, & Q\Theta^2, & Q\Theta^3, & Q\Theta^4 \\ Q, & Q\Theta^3, & Q\Theta^4, & Q\Theta, & Q\Theta^2 \end{pmatrix} \qquad \sigma = \chi.
\]

Letzteres findet man aus den Formeln (1) und (2), wonach
\[
\Theta\chi = \varphi\Theta^3, \quad \Theta^2\chi = \chi\varphi^2\Theta^4, \quad \Theta^3\chi = \varphi^2\Theta, \quad \Theta^4\chi = \varphi\chi\Theta^2.
\]

Man sieht, dass diese Permutationen alle zur ersten Art gehören, und dass also die Permutationsgruppe, um die es sich handelt, keine andere als die alternirende Gruppe von fünf Ziffern (Bd. I, \S.\,177) sein kann. Daraus ergiebt sich auch, dass die Ikosa\"edergruppe einfach ist, und folglich mit der Gruppe übereinstimmt, die wir schon in \S.\,34 vorläufig als Ikosa\"edergruppe bezeichnet haben.

\begin{center}
\S.\,60.\\[2pt]
\gesperrt{Die Grundformen der Ikosa\"edergruppe.}
\end{center}

Wir haben oben die eine der drei Grundformen der Ikosa\"edergruppe $f$ gefunden:
\weq{1}{f = x_1 x_2 (x_1^{10} + 11 x_1^5 x_2^5 - x_2^{10}).}
Es fehlen uns also noch zwei dieser Formen, eine vom 20\textsuperscript{sten} und eine vom 30\textsuperscript{sten} Grade.

Die Grundform 20\textsuperscript{sten} Grades ergiebt sich als die \gesperrt{Hesse}'sche Covariante von $f$. Wir setzen sie
\weq{2}{H = \frac{1}{121}\,[f''(x_1,x_1)\, f''(x_2,x_2) - f''(x_1,x_2)^2],}
%CONT>>
